\section{Related work}
\label{related}

This work draws on three lines of research: prior approaches to coordination in multi-agent learning, the ``neuron as an agent'' paradigm, and Hebbian learning as a mechanism for adaptive behaviour in artificial systems.

\paragraph{Coordination in MARL and LLM-agent systems.} Existing approaches to coordination treat other agents as transient elements of the environment rather than as persistent partners. Three lines of work illustrate this. Model-based auxiliary prediction losses let agents leverage expert behaviour by predicting others' actions \cite{ndousse2021emergent}, but agents minimise the auxiliary objective without forming explicit or lasting relationships, so other agents serve as prediction targets rather than as sources of retained knowledge. Intrinsic motivation methods reward social influence to encourage coordination \cite{jaques2019social}, shaping behaviour without enabling inter-agent skill transfer. Recent Theory of Mind frameworks introduce perspective-taking to improve interpretation of agent intentions \cite{licua2025mindforge}, but enhance reasoning about others rather than learning from them. Across all three, performance gains do not translate into stable social structure \cite{ye2025efficient}.

\paragraph{Neuron as an agent.} A complementary line of research treats individual computational units within a neural network as autonomous, self-interested actors. Early work framed neurons as agents performing Gibbs sampling with synaptic weights modelled as stochastic Gaussian policies \cite{wang2015reinforcement}, and later work used MARL to optimise recurrent networks by viewing neurons as agents that drive the collective system toward rewarded states \cite{chalk2019training}. Other work locate the agent at the synapse rather than the neuron, treating the synapse as the fundamental reinforcement-learning unit responsible for credit assignment \cite{bhargava2022gradient}, or address inter-neuron communication and reward distribution without a trusted third party using auction theory and counterfactual returns \cite{ohsawa2018neuron}. Building on these foundations, \cite{ott2020giving} argues that intelligence emerges bottom-up from the local interactions of self-interested neurons balancing global task rewards against biological incentives. We invert this direction: rather than treating neurons within a network as agents, we treat a population of agents as a network and ask what plasticity rule governs the bonds between them.

\paragraph{Hebbian learning in artificial systems.} Hebbian learning has been used as a local, biologically inspired rule for shaping adaptive behaviour in both neural and multi-agent systems. Prior work leverages Hebbian-style updates to model social factors such as emerging empathy and counter-empathy in agent learning \cite{chen2019reaching}, to perform unsupervised clustering of agents in complex environments such as StarCraft II \cite{kang2023forecasting}, and to build adaptive network models that connect Hebbian plasticity with synchronisation and social bonding dynamics \cite{mukeriia2024multi}. In parallel, Hebbian rules have been explored as a mechanism for robust generalisation, for example by evolving compact shared plasticity rules \cite{pedersen2021evolving} or by meta-learning generalised bidirectional Hebb-style update rules as an alternative to explicit backpropagation \cite{sandler2021meta}. These works apply Hebbian learning either within an agent's internal network or to model behavioural dynamics; we apply it to the inter-agent bond structure itself, treating the population graph as the substrate of plasticity.
